\documentclass[UTF8,a4paper,12pt,fontset=fandol]{ctexart}
\usepackage[left=2.6cm,right=2.4cm,top=2.5cm,bottom=2.5cm]{geometry}
\usepackage{booktabs,longtable,array,ragged2e}
\usepackage{enumitem}
\usepackage{hyperref}
\usepackage{xurl}
\usepackage{fancyhdr}
\usepackage{setspace}
\usepackage{caption}
\usepackage{microtype}
\setCJKmonofont{FandolFang-Regular}
\hypersetup{colorlinks=true,linkcolor=black,urlcolor=blue,citecolor=black,
  pdftitle={基于自我反思与 Critic 反馈的自动驾驶视觉语言动作模型自进化研究}}
\setlength{\parindent}{2em}
\setlength{\parskip}{0.25em}
\setlength{\headheight}{15pt}
\setlength{\LTpre}{0.5em}
\setlength{\LTpost}{0.8em}
\renewcommand{\arraystretch}{1.18}
\setlist{nosep,leftmargin=2.5em}
\pagestyle{fancy}
\fancyhf{}
\fancyhead[C]{自动驾驶 VLA 自进化研究}
\fancyfoot[C]{--- \thepage\ ---}
\renewcommand{\headrulewidth}{0.4pt}
\newcolumntype{P}[1]{>{\RaggedRight\arraybackslash}p{#1}}

\begin{document}
\begin{titlepage}
  \centering
  \vspace*{2.2cm}
  {\zihao{1}\bfseries 基于自我反思与 Critic 反馈的\\[0.35em]
  自动驾驶视觉语言动作模型自进化研究\par}
  \vspace{2.0cm}
  {\zihao{3}课程实践项目报告（选题 4）\par}
  \vspace{0.8cm}
  {\zihao{4}基座方向：OpenDriveVLA / Vision-Language-Action\par}
  {\zihao{4}实验版本：OpenDriveVLA 冻结基座轻量实验 Demo v0.3\par}
  \vfill
  {\zihao{4}2026 年 8 月\par}
\end{titlepage}
\pagenumbering{Roman}
\tableofcontents
\clearpage
\pagenumbering{arabic}
\onehalfspacing
\section*{摘要}
\addcontentsline{toc}{section}{摘要}
现有端到端自动驾驶视觉语言动作模型通常依赖固定数据进行一次性训练，部署后的失败轨迹难以自动转化为新知识。本文提出并实现一个面向课程验证的 Driving Reflection 自进化框架，以本地完整下载并核验的 OpenDriveVLA-0.5B checkpoint 作为冻结基座资产和标准轨迹接口，以“场景输入—轨迹规划—Critic 评价—结构化反思—经验记忆—再训练”为闭环。系统将 Critic 分解为 Safety、Rule、Comfort 三个维度，提供规则评价器、可连接 OpenAI-compatible 服务的 LLM Critic，以及由评价数据拟合的 Reward Model Critic；通过质量过滤构造 5,000 条统一 JSONL 反思样本，并实现 Baseline、SFT、失败样本加权 Reflection SFT、三轮失败回放自进化与轻量 DPO。项目同时提供面向 OpenDriveVLA 的 LoRA SFT、Reflection SFT 和 DPO 数据导出、配置及训练入口，使 Windows 无训练 Demo 与 Linux/\allowbreak{}CUDA 完整训练保持同一数据语义。

在固定随机种子生成的 1,000 条未见合成场景上，Baseline 的碰撞风险代理率、交通违章代理率、行人礼让失败率分别为 8.9\%、45.4\% 和 5.9\%；Reflection SFT 将三项指标降至 3.3\%、7.2\% 和 1.4\%，整体 Critic 得分由 91.176 提升至 97.920；Reflection DPO 进一步达到 3.0\%、5.1\% 和 0.3\%。这些结果验证了“失败定位—纠正知识—再学习”闭环在研究原型上的有效性。本文同时明确：上述结果属于 CPU 可复现的开放环合成代理实验，不等同于 nuScenes 或 CARLA 的真实闭环性能。

关键词：自动驾驶；视觉语言动作模型；自我反思；Critic；监督微调；直接偏好优化；自进化

\section{课题要求核对与研究问题}
\subsection{原始课题要求}
课程文档要求完成四个核心模块：第一，构建 Trajectory→Critic→Reflection 自动生成流水线，形成约 5K～10K 条可过滤、可评价的数据；第二，实现 Safety Score、Rule Score、Comfort Score，并比较 Rule-based、LLM Critic、Reward Model；第三，以 OpenDriveVLA 为基座方向，对比 Baseline、SFT、Reflection SFT，并选做 Reflection+DPO/\allowbreak{}GRPO；第四，在复杂和未见场景中报告碰撞、违章、舒适性、失败定位、Critic 一致性与泛化分析。交付物还必须包括数据、训练、评估全套脚本，配置规范和可现场演示的完整流程。

\subsection{研究问题}
本文围绕四个问题展开：

\begin{enumerate}
  \item 多维 Critic 能否稳定发现跟车、红灯、超速、行人礼让和舒适性问题？
  \item 结构化反思能否把评价证据转换为可训练的纠正策略，而非泛泛文本解释？
  \item Reflection SFT 与偏好优化能否相对普通 SFT 进一步降低安全和规则风险？
  \item 轻量原型如何保留与 OpenDriveVLA、nuScenes、CARLA 的数据和训练接口，使现场演示与后续大模型实验兼容？
\end{enumerate}

\section{指定文献与代码研读}
\subsection{六篇指定论文}
OpenDriveVLA 将多视角图像产生的场景、Agent 和地图结构化 token，通过分层视觉—语言对齐接入 Qwen2.5，并采用四阶段训练：特征对齐、驾驶指令微调、Agent–Environment–Ego 交互学习、轨迹规划微调。轨迹被离散为自回归 waypoint token。该工作为本项目提供基座输入/\allowbreak{}输出范式，但其主要结果为 nuScenes 开放环评估，本身不具备部署后持续进化闭环。

UniAD 的关键思想是“以规划为中心”统一检测、跟踪、地图、运动预测、占用和规划，通过 query 接口让上游任务服务最终规划目标。本项目继承这一目标导向原则：Critic 不评价语言是否华丽，而评价轨迹对安全、规则和舒适性的贡献。

DriveAgent-R1 引入主动感知和混合思考：简单场景采用文本推理，复杂场景调用视觉工具；训练由双模式 SFT、强制对比模式 RL、自动模式选择 GRPO 构成。其奖励同时包含动作准确度、格式一致性与工具调用成本。由此，本项目把 Critic 输出设计为严格结构化 JSON，并把可选在线 LLM 评价与默认离线评价解耦，保证可复现性和扩展性。

Alpamayo-R1 使用 Chain of Causation 数据，把“驾驶决策—可观测因果因素—因果推理链”绑定到轨迹，采用 VLM 与流匹配轨迹解码器，并通过 SFT 与 RL 对齐推理和动作。其最重要启示是避免与动作无关的冗长思维链。本项目的 Reflection 模板因此限定为 verdict、root\_causes、evidence、corrective\_strategy、counterfactual\_action 和 confidence 六类字段。

Agent-Driver 将 LLM 作为认知 Agent，引入工具库、常识/\allowbreak{}经验记忆、推理引擎与自我反思。官方实现依赖预缓存 nuScenes 数据和外部 API。本文借鉴其工具化与经验闭环，但将核心 Critic 和 demo 设计为无 API、无 GPU 也能运行，同时保留在线服务适配接口。

DeepSeek-R1 表明纯强化学习可诱发自验证、反思和长程推理，并采用冷启动数据与多阶段训练稳定语言质量。其方法无法直接照搬到连续驾驶控制：驾驶奖励必须同时约束安全、规则和舒适性，并防止奖励投机。因此本项目先用可验证规则产生锚点，再训练 Reward Model，并以安全权重最高的组合得分驱动过滤和偏好优化。

\subsection{五个指定代码仓库}
\begin{center}
\small
\begin{longtable}{@{}P{0.17\textwidth}P{0.25\textwidth}P{0.20\textwidth}P{0.26\textwidth}@{}}
\toprule
\textbf{仓库} & \textbf{关键可复用内容} & \textbf{对本项目的影响} & \textbf{复现限制} \\
\midrule
\endfirsthead
\toprule
\textbf{仓库} & \textbf{关键可复用内容} & \textbf{对本项目的影响} & \textbf{复现限制} \\
\midrule
\endhead
\midrule
\multicolumn{4}{r}{续下页} \\
\endfoot
\bottomrule
\endlastfoot
DriveLM & nuScenes/\allowbreak{}CARLA 上的 Graph VQA、感知—预测—规划逻辑依赖、评估流程 & Reflection 保留证据链与任务因果关系 & 语言数据含非商业许可，完整数据需单独获取 \\
Alpamayo 1.5 & 推理、导航、多相机、VQA 示例；SFT/\allowbreak{}RL recipes & 报告采用“推理与动作一致性”分析维度 & 权重约 22GB，单样本约 24GB 显存且资源受控 \\
Agent-Driver & \path{agentdriver/}、\path{data/}、\path{experiments/}、\path{scripts/} 分层；工具、记忆、自反思 & 工程结构按数据/\allowbreak{}模型/\allowbreak{}评估/\allowbreak{}demo 解耦 & 需要外部 API 与预缓存 nuScenes 数据 \\
DriveAgent-R1 & \path{tool-rl/} 训练代码、SFT 与 Cascaded RL、多轮工具调用 & LLM Critic 使用严格格式并预留在线接口 & 官方评估脚本、数据和权重尚非全部公开；要求 Ubuntu/\allowbreak{}CUDA \\
OpenDriveVLA & \path{drivevla/}、\path{llava/}、\path{projects/}、\path{scripts/}；0.5B 推理与数据准备 & \path{Policy.plan()} 作为可替换基座边界 & 深度依赖 mmcv/\allowbreak{}mmdet3d、PyTorch、GPU 和 nuScenes \\
\end{longtable}
\end{center}
链接核对发现：课程 Word 中 OpenDriveVLA 显示为 arXiv:2503.23463，但内部跳转关系一度指向 2409.12191；前者标题和官方仓库均与 OpenDriveVLA 一致，故本文采用 2503.23463。DriveAgent-R1 的原 GitHub 地址已重定向至 \path{wczheng/DriveAgent-R1}。

\section{方法}
\subsection{总体闭环}
系统工作流为：场景生成或数据集适配器输出统一 Scenario；Planner 生成未来 12 个 waypoint；三维 Critic 评价轨迹并定位失败；Reflection Builder 把失败与证据映射为纠正动作；Quality Filter 过滤低置信或无证据样本；训练器执行 SFT、Reflection SFT 或 DPO；评估器在独立未见场景上重新测量风险。所有中间产物使用 JSON/\allowbreak{}JSONL，因此可把轻量 Planner 替换为 OpenDriveVLA，而无需改写 Critic、数据过滤和分析模块。

\subsection{场景与轨迹表示}
Scenario 包含自车速度、限速、前车距离/\allowbreak{}速度、信号灯、停止线距离、行人距离、道路曲率、导航指令、天气与未见场景标志。Trajectory 包含 12×0.5s 的 \texttt{[纵向位置,\ 横向位置,\ 速度]} waypoint、目标速度、策略名和简短规划依据。训练/\allowbreak{}验证场景使用 clear/\allowbreak{}rain/\allowbreak{}fog，未见测试额外加入 night，并扩大曲率和红灯比例以形成分布偏移。

\subsection{三维 Critic}
规则 Critic 的安全维度检查时间到碰撞 TTC、行人礼让；规则维度检查超速和红灯停车；舒适性维度检查纵向 jerk 与横向加速度代理量。综合得分为：

\texttt{Score\ =\ 0.45\ ×\ Safety\ +\ 0.35\ ×\ Rule\ +\ 0.20\ ×\ Comfort}。

安全权重最高，反映自动驾驶评价的约束优先级。Rule-based Critic 给出确定性标签和数值证据；LLM Critic 支持 OpenAI-compatible \path{/chat/completions} 严格 JSON 调用，默认离线模式使用可重复的有界分歧代理，避免 API 成本破坏一键复现；Reward Model 使用规则评价样本拟合三输出岭回归，验证学习型评价器能否逼近规则锚点。报告中的 LLM 离线结果明确标注为代理，不声称是真实大模型判断。

\subsection{Driving Reflection 模板}
每条 Reflection 包含：是否接受或修订；根因标签；可核验证据；逐条纠正策略；反事实动作；置信度。根因覆盖 unsafe\_following、pedestrian\_yield\_failure、speeding、red\_light\_risk 和 uncomfortable\_motion。质量分结合反思置信度与证据覆盖度，阈值为 0.58。通过该设计，训练数据不是“下次小心”式空泛文本，而是可定位、可验证、可转成偏好对的纠正知识。

\subsection{Self-Evolution 训练}
Baseline 是带噪声且可能漏检信号灯、行人或前车的轻量规划器。SFT 对专家目标速度执行带 L2 正则的最小二乘拟合。Reflection SFT 对 Critic 定位的失败样本按失败数量提高权重，使长尾错误在训练中获得更强梯度。DPO 轻量实验从 \texttt{专家纠正轨迹（chosen）} 与 \texttt{原失败轨迹（rejected）} 构造偏好对，在 Reflection SFT 权重上执行稳定的小步偏好更新。

为体现持续进化而非单次重训，系统进一步实现三轮失败回放。每轮用当前策略重新生成 1,813 条训练轨迹，Critic 定位失败，Reflection 生成纠正知识并写入经验记忆池，再依据失败严重度和综合得分动态提高样本权重。三轮轻量实验中，平均 Overall 从 96.245 提升到 97.238；需要修订的轨迹由 326 条降至 252 条，失败事件由 387 个降至 302 个，累计形成 835 条可追溯反思记忆。该结果属于冻结基座接口上的轻量代理实验，用于验证自进化机制的方向和工程闭环。

完整模型训练代码位于 \path{scripts/train_opendrivevla_peft.py}。其从本地 OpenDriveVLA-0.5B 初始化，冻结视觉塔和原始参数，仅在 Q/\allowbreak{}K/\allowbreak{}V/\allowbreak{}O 注意力投影注入 LoRA；SFT 与 Reflection SFT 使用交叉熵目标，DPO 使用带冻结参考策略 log-ratio 的偏好损失。Windows 可通过 \path{--check-only} 核验模型、代码、数据和配置，Linux/\allowbreak{}CUDA 环境取消该参数即可执行训练。官方仓库尚未发布训练脚本，因此本项目补充了独立 PEFT 入口。

\subsection{基座模型选择与非从零训练说明}
正式研究路线把 OpenDriveVLA-0.5B 作为必要基座，而不是从随机参数训练。该选择直接对应课题指定的 OpenDriveVLA 论文与官方代码，具备驾驶场景 token、Qwen2.5 语言骨干和 waypoint 动作输出，且 0.5B 规模较 7B 或 22GB 级模型更适合课程算力。训练时必须加载其预训练 checkpoint，冻结视觉编码器和大部分语言层，仅对 multimodal projector、trajectory head 与少量 LoRA 参数执行 Reflection SFT，再使用 chosen/\allowbreak{}rejected 轨迹对进行 DPO；配合 8-bit/\allowbreak{}4-bit 量化、缓存视觉特征与梯度累积，可进一步降低显存需求。

本地已完整下载 OpenDriveVLA-0.5B 权重与官方代码。审计脚本验证 \path{model.safetensors} 大小为 1,466,080,342 字节、Safetensors 文件头有效、模型架构为 \path{LlavaQwenForCausalLM}，且 \path{drivevla/}、\path{llava/}、\path{projects/}、\path{scripts/}、\path{third_party/} 均存在。本地 Demo 采用冻结基座适配模式：\path{BaseModelAdapter} 固定记录 OpenDriveVLA checkpoint、架构、tokenizer 和官方六点轨迹格式；若存在 Linux/\allowbreak{}GPU 产生的匹配缓存，则回放真实 OpenDriveVLA 输出，否则明确使用 \texttt{LiteVLA\ CPU\ surrogate} 完成轻量实验。页面分别披露 \path{checkpoint_installed}、\path{weights_loaded} 和轨迹来源，避免把“模型已下载”误写成“本机已执行视觉推理”。

\section{实验设计}
\subsection{数据规模与划分}
固定 \path{seed=42} 生成 5,000 条样本。其中训练/\allowbreak{}验证来自常见天气，最后 1,000 条为未见测试场景并包含夜间条件。流水线按轮换方式使用 80\% Rule Critic、10\% LLM 离线代理、10\% Reward Model；经过质量阈值过滤保留 2,730 条，最终 Reflection 训练使用 1,813 条训练划分高质量样本。

为与完整 OpenDriveVLA 训练对齐，数据导出器进一步生成 1,813 条普通 SFT 样本、1,813 条 Reflection SFT 样本和 1,112 对 DPO chosen/\allowbreak{}rejected 偏好数据。每条记录包含基座名称、场景 ID、来源类型、结构化 prompt、轨迹 token 和预留的多相机 \path{image_refs} 字段；轻量实验中图像引用为空，接入 nuScenes 时保持 schema 不变并替换为真实图像和 UniAD 特征引用。

\subsection{指标}
主要指标包括：碰撞风险代理率（TTC 失败）、交通违章代理率（超速和红灯风险）、行人礼让失败率、不舒适率、Safety/\allowbreak{}Rule/\allowbreak{}Comfort/\allowbreak{}Overall 均分，以及目标速度相对专家规则的 MAE。由于轻量模拟器没有真实物理碰撞和传感器闭环，所有“率”均明确称为代理指标。

\subsection{对比与消融}
对比四种策略：Baseline、SFT、Reflection SFT、Reflection DPO。Critic 分析比较 Rule、LLM 离线代理与 Reward Model 的样本数、均值和标准差；失败案例分析聚焦红灯、行人和前车三类风险；泛化测试仅使用训练中未出现的 night 天气和更大曲率范围。

\section{实验结果与分析}
\subsection{未见场景主结果}
\begin{center}
\small
\begin{longtable}{@{}P{0.14\textwidth}P{0.09\textwidth}P{0.09\textwidth}P{0.09\textwidth}P{0.09\textwidth}P{0.08\textwidth}P{0.10\textwidth}P{0.09\textwidth}@{}}
\toprule
\textbf{策略} & \textbf{碰撞风险率\ensuremath{\downarrow}} & \textbf{违章率\ensuremath{\downarrow}} & \textbf{行人礼让失败率\ensuremath{\downarrow}} & \textbf{Safety\ensuremath{\uparrow}} & \textbf{Rule\ensuremath{\uparrow}} & \textbf{Overall\ensuremath{\uparrow}} & \textbf{目标速度 MAE\ensuremath{\downarrow}} \\
\midrule
\endfirsthead
\toprule
\textbf{策略} & \textbf{碰撞风险率\ensuremath{\downarrow}} & \textbf{违章率\ensuremath{\downarrow}} & \textbf{行人礼让失败率\ensuremath{\downarrow}} & \textbf{Safety\ensuremath{\uparrow}} & \textbf{Rule\ensuremath{\uparrow}} & \textbf{Overall\ensuremath{\uparrow}} & \textbf{目标速度 MAE\ensuremath{\downarrow}} \\
\midrule
\endhead
\midrule
\multicolumn{8}{r}{续下页} \\
\endfoot
\bottomrule
\endlastfoot
Baseline & 8.9\% & 45.4\% & 5.9\% & 94.909 & 81.335 & 91.176 & 5.163 \\
SFT & 3.2\% & 7.5\% & 1.5\% & 98.339 & 95.985 & 97.847 & 1.512 \\
Reflection SFT & 3.3\% & 7.2\% & 1.4\% & 98.394 & 96.123 & 97.920 & 1.561 \\
Reflection DPO & 3.0\% & 5.1\% & 0.3\% & 99.045 & 97.343 & 98.640 & 1.980 \\
\end{longtable}
\end{center}
SFT 带来最大的第一阶段提升，说明专家动作监督仍是基本盘。Reflection SFT 相比 SFT 的总体提升较小，但 Rule 分数、违章率和行人礼让均继续改善，符合失败样本加权主要强化安全长尾而非平均回归误差的预期。DPO 进一步显著降低行人和规则风险，但目标速度 MAE 从 1.561 上升到 1.980，表现出更保守的安全偏好与专家平均速度之间的权衡；这正是偏好优化需要结合效率约束的原因。

\subsection{Critic 对比}
\begin{center}
\small
\begin{longtable}{@{}P{0.14\textwidth}P{0.12\textwidth}P{0.15\textwidth}P{0.14\textwidth}P{0.31\textwidth}@{}}
\toprule
\textbf{Critic} & \textbf{样本数} & \textbf{Overall 均值} & \textbf{标准差} & \textbf{特点} \\
\midrule
\endfirsthead
\toprule
\textbf{Critic} & \textbf{样本数} & \textbf{Overall 均值} & \textbf{标准差} & \textbf{特点} \\
\midrule
\endhead
\midrule
\multicolumn{5}{r}{续下页} \\
\endfoot
\bottomrule
\endlastfoot
Rule-based & 4000 & 92.132 & 11.632 & 确定、可解释、覆盖受规则限制 \\
LLM offline proxy & 500 & 91.959 & 11.986 & 可复现的接口代理，不等同真实 LLM \\
Reward Model & 500 & 91.213 & 8.374 & 平滑但方差偏小，可能弱化极端失败 \\
\end{longtable}
\end{center}
三个 Critic 均值接近，但 Reward Model 的标准差明显更小，说明线性奖励模型存在回归均值倾向，对严重长尾风险可能不够敏感。因此推荐采用“规则硬门槛 + 学习型排序 + LLM 解释复核”的组合，而非让单一 Reward Model 决定安全样本是否通过。

\subsection{失败归因与泛化}
红灯失败可由“红灯、停止线距离、目标速度过高”三项证据定位，纠正策略为提前渐进制动并保持停车；行人失败关联行人距离与当前目标速度，纠正为提前礼让并保留全停余量；跟车失败使用 TTC 代理直接定位。未见场景仍保持显著改善，说明闭环学习到的是条件约束而非场景编号。然而，夜间在本 demo 中只是结构化变量，没有真实图像域偏移，因此不能据此宣称视觉泛化已经解决。

\subsection{舒适性指标解释}
四种策略在当前轨迹生成器中均未触发不舒适阈值，Comfort 均分为 100。这不是“舒适性完美”，而是当前解析轨迹由平滑一阶速度控制器生成，场景难度不足以区分模型。真实扩展应加入急切入、非平稳前车、道路颠簸和横向避障，并使用加速度、jerk、横摆角速度及 CARLA 动力学反馈扩大舒适性辨识度。

\section{Demo 与工程实现}
运行 \texttt{python\ scripts/\allowbreak{}run\_all.py\ --samples\ 5000} 可一键执行数据构建、轻量训练和评估；随后执行 \texttt{python\ scripts/\allowbreak{}prepare\_vla\_training.py} 导出 OpenDriveVLA 三阶段训练数据。运行 \texttt{python\ scripts/\allowbreak{}run\_demo.py} 后访问 \path{http://127.0.0.1:8000}，Windows 也可双击根目录的 \texttt{启动Demo.bat}。新版网页由 Python 动态 API 驱动，不是离线静态页面：前端不保存固定场景、轨迹或分数，四类演示场景由后端从完整 5,000 条训练数据中动态匹配，也支持按 \path{scene_id} 和随机偏移加载。道路几何、停止线、信号灯、前车、行人、弯道和天气效果均由本次场景 JSON 驱动；红色轨迹表示初始 Baseline，青色轨迹表示反思后策略；播放控件动态模拟 6 秒路径规划窗口。

现场演示建议按以下顺序操作：

\begin{itemize}
  \item 选择“红灯路口”，运行后观察 Baseline 是否因较高目标速度触发 \path{red_light_risk}。
  \item 播放双轨迹动画，对比反思后策略的提前减速和 Overall 变化。
  \item 在时间线中依次解释感知、初始规划、Critic 证据、纠正策略和重规划结果。
  \item 切换“行人横穿”或“前车急缓”，验证同一闭环对不同失败根因的迁移。
  \item 切换 Reflection SFT 与 Reflection DPO，讨论偏好优化更安全但可能更保守的权衡。
\end{itemize}

页面可继续手工调整速度、距离、曲率、导航和天气；修改后使用新的场景 ID，避免误命中旧缓存。服务端使用当前策略实时规划，并以规则 Critic 65\%、训练后的 Reward Critic 25\%、完整训练集七近邻校准 10\% 计算三维分数。\path{/api/logs} 持续输出数据文件加载、Reward Critic 参数连接、规划调用的实际 runtime 与耗时、评分近邻和 Reflection 结果；\path{/api/data/status}、\path{/api/scenarios}、\path{/api/scenario} 提供可核验的数据连接。完整操作、字段含义和故障排查见 \texttt{docs/\allowbreak{}Demo使用说明.md}；基座选择与轻量化边界见 \texttt{docs/\allowbreak{}基座模型与轻量化说明.md}。

主要目录为：\path{configs/} 保存实验配置；\path{src/selfevolve_drive/} 保存 schema、模拟器、Planner、Critic、Reflection、训练与评估；\path{scripts/} 保存独立数据/\allowbreak{}训练/\allowbreak{}评估/\allowbreak{}demo 入口；\path{data/} 保存 5K JSONL；\path{models/} 保存 SFT、Reflection SFT、DPO 与 Reward Model 参数；\path{outputs/} 保存真实运行指标；\path{tests/} 保存回归测试；\path{references/} 保存资料清单；\path{report/} 保存报告。

\section{OpenDriveVLA /\allowbreak{} nuScenes /\allowbreak{} CARLA 扩展方案}
第一阶段接入 nuScenes：使用 OpenDriveVLA 官方数据准备脚本生成多视角视觉 token、ego state、command 和专家 waypoint；使用现有 \path{OpenDriveVLAAdapter} 与官方轨迹解析器返回本项目 Trajectory schema；用 nuScenes box、地图与信号标注替换合成 Scenario。第二阶段使用已交付的 \path{opendrivevla_peft.json} 和 \path{train_opendrivevla_peft.py} 执行 Baseline、普通轨迹 LoRA SFT、\texttt{输入+失败轨迹+Reflection→纠正轨迹} Reflection SFT，并从同场景候选轨迹构造标准 DPO 对。第三阶段接入 CARLA/\allowbreak{}Bench2Drive 闭环，测量实际碰撞率、红灯/\allowbreak{}停车牌违章、路线完成率、加速度、jerk 和闭环驾驶分数。

LLM Critic 的线上实验需固定模型版本、temperature=0、JSON schema、重试策略和原始响应缓存；至少抽样 200 条由两名人工标注者复核，报告安全/\allowbreak{}规则/\allowbreak{}舒适三维相关系数、严重失败召回率、Cohen's \ensuremath{\kappa} 或 Krippendorff's \ensuremath{\alpha}。Reward Model 应从线性基线升级为能处理视觉与轨迹 token 的模型，并保留规则硬约束作为安全兜底。

\section{局限、风险与伦理}
本项目已下载并核验 gated OpenDriveVLA-0.5B checkpoint 和官方代码，但尚未下载 nuScenes 完整数据，也未在真实车辆或 CARLA 闭环中训练。Windows 轻量 Planner 不处理原始多相机图像，因而当前结果是以冻结基座资产、标准轨迹格式和统一适配接口为约束的轻量自进化实验，而非官方 OpenDriveVLA 视觉性能复现。LLM Critic 可能产生貌似合理但错误的解释，Reward Model 可能被策略钻空子，自训练还可能放大自身偏差。因此任何自动反思样本都必须经过规则门槛、置信过滤、分布监控与人工抽检；真实车辆部署必须保留独立安全层、故障回退和版本回滚，禁止未经闭环仿真和安全审查直接上车。

\section{结论}
本文完成了课程要求的四大核心模块、全套脚本、5K 数据、三类 Critic、四种策略、三轮持续自进化、未见场景评估和可交互 Demo。轻量实验证明，结构化 Reflection 可以把失败证据转为具有训练价值的纠正信号；普通 SFT 提供主要能力提升，Reflection 加权、失败经验回放与 DPO 对安全长尾有额外收益。项目以已下载并核验的 OpenDriveVLA-0.5B 为冻结基座资产，交付了 LoRA SFT、Reflection SFT、DPO 的数据、配置、环境校验和训练入口；Windows 现场版本通过明确披露的 CPU 回退保障稳定运行。后续只需在 Linux/\allowbreak{}CUDA 和 nuScenes 条件下执行同一训练入口、导入真实轨迹缓存并补充官方指标，无需重写 Critic、Reflection、数据 schema 或 Demo。

\section*{参考文献与代码}
\addcontentsline{toc}{section}{参考文献与代码}
[1] Zhou X. et al. OpenDriveVLA: Towards End-to-end Autonomous Driving with Large Vision Language Action Model. arXiv:2503.23463, 2025. \url{https://arxiv.org/abs/2503.23463}\par
[2] Hu Y. et al. Planning-oriented Autonomous Driving (UniAD). arXiv:2212.10156, 2022. \url{https://arxiv.org/abs/2212.10156}\par
[3] Zheng W. et al. DriveAgent-R1: Advancing VLM-based Autonomous Driving with Active Perception and Hybrid Thinking. arXiv:2507.20879, 2025. \url{https://arxiv.org/abs/2507.20879}\par
[4] NVIDIA. Alpamayo-R1: Bridging Reasoning and Action Prediction for Generalizable Autonomous Driving in the Long Tail. arXiv:2511.00088, 2025. \url{https://arxiv.org/abs/2511.00088}\par
[5] Mao J. et al. A Language Agent for Autonomous Driving. arXiv:2311.10813, 2023. \url{https://arxiv.org/abs/2311.10813}\par
[6] DeepSeek-AI. DeepSeek-R1: Incentivizing Reasoning Capability in LLMs via Reinforcement Learning. arXiv:2501.12948, 2025. \url{https://arxiv.org/abs/2501.12948}\par
[7] OpenDriveLab. DriveLM. \url{https://github.com/OpenDriveLab/DriveLM}\par
[8] NVIDIA. Alpamayo 1.5. \url{https://github.com/NVlabs/alpamayo1.5}\par
[9] Physical Superintelligence Lab. Agent-Driver. \url{https://github.com/physical-superintelligence-lab/Agent-Driver}\par
[10] Zheng W. DriveAgent-R1 code. \url{https://github.com/wczheng/DriveAgent-R1}\par
[11] DriveVLA. OpenDriveVLA code. \url{https://github.com/DriveVLA/OpenDriveVLA}\par
\appendix
\section{交付验收矩阵}
\begin{center}
\small
\begin{longtable}{@{}P{0.22\textwidth}P{0.49\textwidth}P{0.17\textwidth}@{}}
\toprule
\textbf{要求} & \textbf{实现证据} & \textbf{状态} \\
\midrule
\endfirsthead
\toprule
\textbf{要求} & \textbf{实现证据} & \textbf{状态} \\
\midrule
\endhead
\midrule
\multicolumn{3}{r}{续下页} \\
\endfoot
\bottomrule
\endlastfoot
5K～10K 反思数据 & \path{data/reflection_dataset.jsonl}，实际 5,000 条 & 完成 \\
自动生成与质量过滤 & \path{pipeline.py}、质量阈值与 accepted 字段 & 完成 \\
Safety/\allowbreak{}Rule/\allowbreak{}Comfort & \path{critics.py} 三维评分及证据 & 完成 \\
Rule/\allowbreak{}LLM/\allowbreak{}Reward Model & 三类 Critic 与统计对比 & 完成；LLM 默认离线代理 \\
Baseline/\allowbreak{}SFT/\allowbreak{}Reflection SFT & 三组模型及未见场景结果 & 完成 \\
Reflection+DPO/\allowbreak{}GRPO & DPO 轻量实现；GRPO 方案说明 & DPO 完成，GRPO 为扩展 \\
碰撞/\allowbreak{}违章/\allowbreak{}舒适评估 & \path{outputs/metrics.json} & 完成（合成代理） \\
失败原因定位 & 五类 failure 标签与 evidence & 完成 \\
泛化能力 & 1,000 条含夜间/\allowbreak{}高曲率未见场景 & 完成（结构化代理） \\
全套脚本和配置 & \path{scripts/}、\path{configs/default.json} & 完成 \\
一键复现 & \path{scripts/run_all.py} & 完成 \\
现场 Demo & \path{scripts/run_demo.py} + 网页 & 完成 \\
真实场景可视化 & 双策略动态轨迹、信号灯/\allowbreak{}行人/\allowbreak{}前车、反思时间线 & 完成（开放环模拟） \\
必要基座 & OpenDriveVLA-0.5B checkpoint 与官方代码已下载并通过结构审计 & 完成（冻结基座资产） \\
多轮自进化 & 三轮 rollout→Critic→Reflection→记忆回放；835 条经验记忆 & 完成（轻量实验） \\
OpenDriveVLA 训练数据 & SFT 1813、Reflection SFT 1813、DPO 1112 对 & 完成（轻量数据） \\
LoRA/\allowbreak{}PEFT 训练代码 & \path{train_opendrivevla_peft.py} 与配置；Windows 可 check-only & 完成代码；GPU 执行为扩展 \\
使用文档 & \texttt{docs/\allowbreak{}Demo使用说明.md}、\texttt{docs/\allowbreak{}基座模型与轻量化说明.md} & 完成 \\
自动测试 & \path{tests/test_core.py} & 完成 \\
\end{longtable}
\end{center}
\end{document}
